\section{NFL Performance Metrics}
\label{sec:nfl-metrics}

\subsection{Expected Points Added}
\label{sec:epa-framework}

Expected Points Added (EPA) is the most widely used analytical metric in contemporary NFL discourse. The framework assigns a point value to every game state  -- defined by down, distance, yard line, and time remaining  -- representing the expected number of points the offense will score on the current drive from that state. EPA for a given play is then the difference between the expected points after the play and the expected points before the play: a play that moves the team from a state worth 1.5 expected points to a state worth 3.0 expected points generates +1.5 EPA. Conversely, a turnover that transfers possession generates strongly negative EPA for the offense because the post-play state shifts from the offense's expected scoring opportunity to the opponent's.

The conceptual framework traces to the work of \textcite{carter1971operations}, who published one of the first expected point models based on field position and down. The modern implementation, operationalized through the \texttt{nflfastR} package \parencite{carl2021nflfastr} and building on Brian Burke's public work in the late 2000s, employs a multinomial outcome model trained on historical play-by-play data to estimate expected points for every possible game state. \textcite{yurko2019nflwar} provided the most rigorous peer-reviewed description of the EPA framework's application to player evaluation, demonstrating its utility for decomposing team wins into individual player contributions.

EPA's dominance in public NFL analytics reflects genuine strengths. It provides a common currency  -- points  -- for comparing plays of different types (runs versus passes, first downs versus turnovers), enabling apples-to-apples evaluation across the full spectrum of football actions. It is contextual: a 5-yard gain on 3rd-and-4 is correctly valued as substantially more beneficial than the same 5-yard gain on 3rd-and-15, because the former results in a first down (a large positive state transition) while the latter leaves the team in a punting situation. And it is publicly available through \texttt{nflfastR}, enabling independent analysis without proprietary barriers.

However, EPA has several limitations that become acute when it is used to evaluate team quality rather than individual plays. These limitations motivate the drive-level approach adopted in this paper.

\subsubsection{Play-Level Noise and Single-Play Dominance}

American football plays exhibit enormous variance in outcome. A single 80-yard touchdown pass generates EPA equivalent to perhaps 20 short completions combined. When EPA is aggregated to the game or season level to measure team quality, these high-variance events can dominate the resulting totals. A team's season EPA can be disproportionately influenced by a handful of explosive plays that may reflect scheme, talent, or random fortune in roughly equal measure. This noise-to-signal ratio is particularly problematic for cross-team comparison: distinguishing the 5th-best offense from the 15th-best offense on the basis of season EPA totals requires separating a modest quality difference from the cumulative variance of approximately 1,100 individual plays per team per season.

Drive-level aggregation  -- as employed by the \xscore{} model  -- largely mitigates this problem. Rather than summing the point values of individual plays, the drive-level approach classifies each drive into one of four outcome categories (touchdown, field goal, turnover, or punt/other) and estimates the probability of each outcome given the situational state at drive start. The categorical outcome at the drive level is inherently less noisy than the continuous EPA sum, because a single explosive play within a drive contributes only one outcome observation rather than an outsized continuous value to a summation.

\subsubsection{Arbitrary Value Allocation Between Plays}

A fundamental conceptual issue with play-level EPA is the attribution of value to individual plays within a sequence. Consider a drive that begins at the offense's own 25-yard line and scores a touchdown through ten consecutive plays. Each play receives EPA equal to the state-value difference it created. But the ``true'' value of, say, the seventh play in the sequence depends entirely on what came before (the first six plays created the opportunity) and what comes after (the final three plays complete the scoring drive). EPA treats each play as an independent value-creating event, when in reality football drives are cooperative sequences where the marginal contribution of each play depends on the full trajectory.

This attribution problem is analogous to the Shapley value allocation problem in cooperative game theory \parencite{lundberg2017shap}: assigning individual contributions to members of a team whose output is jointly determined. At the drive level, this problem disappears, because the unit of analysis encompasses the full cooperative sequence. The \gds{} framework does not ask which play within a drive created value; it asks whether the entire drive produced a better or worse outcome than expected given its starting conditions.

\subsubsection{Context Sensitivity and Game Script Effects}

EPA does not explicitly model the effect of game script on play-calling and therefore on statistical accumulation. Teams protecting a large lead adopt conservative play-calling: running the ball to consume clock, avoiding deep passes that risk interceptions, and accepting short gains that maintain possession. This game-script effect systematically suppresses EPA per play for teams that frequently hold large leads  -- precisely the best teams in the league. Conversely, teams trailing late must throw aggressively, generating both high-EPA completions and high-negative-EPA turnovers. Season-level EPA totals therefore conflate team quality with game-script exposure, creating a compression effect where the best teams' offensive EPA is artificially suppressed by their tendency to protect leads, and the worst teams' offensive EPA is artificially inflated by their tendency to throw frequently in desperate situations.

The \xscore{} model addresses this directly through its inclusion of score differential as an input feature (\Cref{sec:xscore}). By conditioning expected drive-outcome probabilities on the current scoring margin, the model's baseline already accounts for the reduced scoring expectations when a team is protecting a large lead. A conservative drive that fails to score when the team leads by 28 points generates minimal negative \xvoa{} because the baseline expectation was already low. This game-state conditioning represents a fundamental advantage of the \gds{} approach over raw EPA aggregation for team quality measurement.

\subsection{Defense-Adjusted Value Over Average}
\label{sec:dvoa}

Football Outsiders' Defense-adjusted Value Over Average (DVOA) has been the most widely cited team quality metric in serious NFL analysis for nearly two decades \parencite{schatz2005dvoa}. Like the \gds{} framework, DVOA decomposes team performance into offensive, defensive, and special teams components expressed as percentages relative to league average, enabling analysts to identify not merely that a team is good or bad, but specifically why. A team with +20\% offensive DVOA and $-5$\% defensive DVOA (where negative is better for defense) is identified as an offensive powerhouse with a below-average defense, a characterization that informs strategic analysis in ways that a single composite rating cannot.

DVOA operates at the play level, applying an EPA-like value framework to each down, and then adjusts for two factors that raw EPA ignores: opponent quality and game situation. The opponent adjustment iteratively recalibrates each play's value based on the quality of the opposing unit, so that a play achieving +0.3 EPA against the league's best defense is valued more highly than the same +0.3 EPA against the worst defense. The situational adjustment accounts for second-half performance when teams are trailing or leading by large margins, attempting to down-weight garbage-time performance that does not reflect competitive quality.

\subsubsection{The Reproducibility Problem}

Despite its analytical sophistication, DVOA suffers from a fatal limitation for academic use: it is proprietary and non-reproducible. The complete methodology has never been published in sufficient detail for independent replication. Key aspects remain undisclosed: the precise opponent-adjustment algorithm and its convergence properties, the weighting function applied to garbage-time plays, the normalization procedure that converts raw play values to the DVOA percentage scale, and the handling of special situations such as kneel-downs, spikes, and penalties. No external researcher has ever independently replicated DVOA's published figures, because the information required to do so has never been made available.

This opacity creates a fundamental epistemological problem. When DVOA identifies a team as the league's best defense, the claim rests entirely on the authority of Football Outsiders rather than on verifiable methodology. Researchers citing DVOA must treat it as a black box: they can report its outputs but cannot assess whether those outputs would survive alternative methodological choices. For a question as consequential as ``should an NFL franchise invest \$100 million in defensive talent?''  -- a decision that the ``defense wins championships'' hypothesis directly informs  -- relying on an unverifiable metric is scientifically unsatisfying. This absence of reproducible team quality metrics hampers rigorous evaluation of roster-construction strategies: researchers studying questions such as draft value \parencite{massey2013loserscurse} or salary cap allocation \parencite{mulholland2019optimizing} must either rely on opaque proprietary outputs or construct their own measures from scratch.

The \gds{} framework developed in this paper addresses this gap directly. Every component of the analytical pipeline  -- from feature engineering through model training, calibration, decomposition, and hypothesis testing  -- is fully specified in \Cref{ch:methodology} and implemented in publicly available code. The commitment to transparency is not merely an academic virtue but a practical necessity when the downstream decisions involve hundreds of millions in annual cap spending.

\subsection{Pro Football Focus Grades}
\label{sec:pff}

Pro Football Focus (PFF) employs trained analysts to grade every player on every play of every NFL game, producing player-level grades on a $0$--$100$ scale that can be aggregated to position groups and teams. PFF grades have become deeply embedded in NFL discourse \parencite{pff2014grades}: media coverage routinely cites them, front offices reportedly subscribe to PFF's analytical services, and fan communities treat PFF grades as authoritative assessments of player quality.

PFF's approach represents a fundamentally different epistemological tradition from the model-based metrics discussed above. Where EPA and DVOA derive value from statistical models applied to observable outcomes (yards gained, first downs converted, touchdowns scored), PFF grades derive from human expert judgment applied to process rather than outcome. An analyst watching film grades a quarterback not on whether the pass was completed, but on whether the read was correct, the mechanics were sound, and the ball placement was optimal given the coverage. This process-over-outcome philosophy has genuine merit: it can capture quality that outcome-based metrics miss (a perfectly thrown ball that a receiver drops, or a well-disguised coverage that the quarterback correctly identifies and avoids).

However, PFF grades compound the reproducibility problem identified with DVOA. Not only is the weighting scheme that converts play-level grades to season-level player ratings proprietary, but the underlying data itself  -- the play-by-play grades  -- is generated through a process that is inherently non-reproducible. Two trained analysts watching the same play may assign different grades based on differing interpretations of assignment, effort, and technique. The same analyst may grade differently on different days depending on fatigue, context, and anchoring effects. There is no way for an external researcher to verify that a given player grade is ``correct'' because there is no objective ground truth against which subjective assessments can be validated.

For the purposes of this paper, PFF grades are methodologically unsuitable as a basis for hypothesis testing. The ``defense wins championships'' question requires team quality decomposition into offensive and defensive components, measured on a common scale that permits direct comparison. PFF's subjective, proprietary, and non-reproducible approach fails all three requirements: the scale is anchored to analyst judgment rather than observed outcomes, the methodology cannot be replicated, and the grading process introduces unquantifiable measurement error that cannot be modeled or controlled for.

\subsection{Win Probability Models}
\label{sec:win-probability}

Win probability models estimate the likelihood that a given team will win the game conditional on the current game state, updating after every play to reflect the evolving competitive situation \parencite{stern1991probability}. These models have become ubiquitous in NFL broadcasting (the familiar ``Win Probability'' graphic displayed during telecasts) and in analytical applications such as fourth-down decision-making and late-game strategy evaluation.

\textcite{lock2014randomforests} developed one of the most rigorous academic treatments of NFL win probability, employing random forests to predict game outcomes from pre-play game states. Their model demonstrated that tree-based ensemble methods could capture the complex non-linear relationships between game-state variables (score, time, possession, field position) and win probability more accurately than parametric logit models previously used for this purpose. Win probability models share the \xscore{} model's commitment to situational conditioning  -- both ask ``what would we expect to happen given the current state?''  -- but differ in their target variable. Win probability predicts the binary game outcome (win/loss); \xscore{} predicts the multinomial drive outcome distribution (touchdown, field goal, turnover, or punt). The drive-level target is chosen specifically to enable the three-component decomposition that win probability models cannot provide: a win probability model attributes value to the team as a whole, without distinguishing whether the offense, defense, or special teams is responsible for a favorable game state.

Win probability models are therefore complementary to rather than competitive with the \gds{} framework. They answer the question ``how likely is this team to win from here?'' while \gds{} answers ``how much better than average is this team's offense (or defense) performing?'' The first question is useful for in-game decision-making; the second is necessary for the strategic resource allocation analysis that motivates this paper.
